\documentclass[9pt,a4paper]{extarticle}

% --- Packages ---
\usepackage[utf8]{inputenc}
\usepackage[T1]{fontenc}
\usepackage{lmodern}
\usepackage[spanish,es-tabla]{babel}
\usepackage[a4paper,margin=0.30in,top=0.24in,bottom=0.22in]{geometry}
\usepackage{multicol}
\usepackage{enumitem}
\usepackage{xcolor}
\usepackage{titlesec}
\usepackage{booktabs}
\usepackage{tabularx}
\usepackage{tcolorbox}
\usepackage{hyperref}
\usepackage{amsmath,amssymb}
\usepackage{microtype}
\usepackage{tikz}
\usetikzlibrary{arrows.meta,positioning,calc,fit,backgrounds,shadows}

\definecolor{boxbg}{gray}{0.96}
\definecolor{boxframe}{gray}{0.15}
\definecolor{soft}{gray}{0.86}
\definecolor{cuno}{RGB}{0,98,168}   % ciclo 1: turno a turno
\definecolor{cdos}{RGB}{214,110,0}  % ciclo 2: revisión
\definecolor{ctres}{RGB}{190,30,45} % ciclo 3: callejón
\definecolor{cpar}{RGB}{0,128,90}   % parada
\newcommand{\sw}[1]{\textcolor{#1}{\rule{1.6ex}{1.6ex}}}

\hypersetup{colorlinks=true,linkcolor=black,citecolor=black,urlcolor=black,
    pdftitle={IA Generativa - Entrega 3: Sistema de agentes para turnos medicos},pdfauthor={Grupo 11}}

\titleformat{\section}{\fontsize{8.6}{9.6}\bfseries\scshape}{}{0em}{}[\vspace{-3.6pt}\rule{\linewidth}{0.55pt}\vspace{-3.2pt}]
\titlespacing*{\section}{0pt}{2.2pt}{1.5pt}

\setlength{\columnsep}{11pt}
\setlength{\parskip}{1pt}
\setlength{\parindent}{0pt}
\setlist{noitemsep,topsep=0.6pt,partopsep=0pt,parsep=0.1pt,leftmargin=0.95em}
\pagestyle{empty}
\newcommand{\tabfont}{\fontsize{6.7}{7.7}\selectfont\setlength{\tabcolsep}{2.6pt}}
\newcommand{\notefont}{\fontsize{6.2}{7.1}\selectfont}
\newcolumntype{L}[1]{>{\raggedright\arraybackslash}p{#1}}
\newcolumntype{Y}{>{\raggedright\arraybackslash}X}

\begin{document}
\fontsize{7.2}{8.2}\selectfont

% ============================ HEADER ============================
\begin{tcolorbox}[colback=boxbg,colframe=boxframe,arc=1.5mm,boxrule=0.8pt,top=1.5pt,bottom=1.5pt,left=6pt,right=6pt]
\begin{minipage}[c]{0.62\textwidth}
    {\fontsize{11}{12.5}\selectfont \textbf{Asignación de Turnos Médicos con un Sistema de Agentes}}\\[1pt]
    {\fontsize{8}{9}\selectfont \textbf{Deliverable 3 \textbar\ Generative Artificial Intelligence (580694)}}
\end{minipage}\hfill
\begin{minipage}[c]{0.37\textwidth}\raggedleft
    {\fontsize{7.6}{8.8}\selectfont \textbf{Universidad de Concepción \textbar\ Grupo 11}}\\
    {\fontsize{6.6}{7.8}\selectfont \textbf{Repo:} \href{https://github.com/matirocha/genai-grupo11}{\texttt{github.com/matirocha/genai-grupo11}}}\\
    {\fontsize{6.6}{7.8}\selectfont \textbf{Video:} \href{VIDEO_LINK}{\texttt{VIDEO\_LINK}}}
\end{minipage}
\end{tcolorbox}

\vspace{-5pt}
% ============================ AGENT GRAPH ============================
\begin{center}
\begin{tikzpicture}[x=1mm,y=1mm,
    font=\fontsize{6}{6.9}\selectfont,
    tool/.style={draw=boxframe,sharp corners,fill=boxbg,line width=0.4pt,align=center,inner sep=1.4pt,minimum height=7.5mm,
                 execute at begin node={\hyphenpenalty 10000\relax}},
    agent/.style={draw=boxframe,rounded corners=3mm,fill=soft,line width=1.4pt,align=center,inner sep=1.8pt,inner xsep=2.6pt,
                  minimum height=9mm,drop shadow={opacity=0.4,shadow xshift=0.6mm,shadow yshift=-0.6mm},
                  execute at begin node={\hyphenpenalty 10000\relax}},
    lbl/.style={font=\fontsize{5.5}{6.3}\selectfont,align=center,fill=white,inner sep=0.5pt},
    cyc/.style={font=\fontsize{5.6}{6.4}\selectfont\bfseries},
    arr/.style={-{Latex[length=1.3mm]},line width=0.5pt},
    c1/.style={draw=cuno,fill=cuno!10},  c1a/.style={draw=cuno,fill=cuno!35},
    c2/.style={draw=cdos,fill=cdos!10},  c2a/.style={draw=cdos,fill=cdos!35},
    c3/.style={draw=ctres,fill=ctres!10}, c3a/.style={draw=ctres,fill=ctres!35},
    cp/.style={draw=cpar,fill=cpar!12}]
  % carril normal
  \node[tool,text width=10mm] (in) at (0,0) {Instancia\\JSON};
  \node[tool,text width=23mm,line width=0.8pt] (orq) at (23,0) {\textbf{Coordinador de turnos}\\código, no agente:\\estado y turno $k$};
  \node[tool,c1,text width=19mm] (t1) at (55,0) {\textbf{T1 opciones}\\\textbf{viables}\\filtro legal (D2)};
  \node[agent,c1a,text width=19mm] (a1) at (91,0) {\textbf{A1 Asignador}\\elige 1 opción\\del menú};
  \node[tool,c1,text width=17mm] (t2) at (120,11) {\textbf{T2 perfilar}\\alertas R1--R3};
  \node[agent,c2a,text width=18mm] (a2) at (149,0) {\textbf{A2 Crítico}\\acepta, rechaza\\o sondea};
  \node[tool,c2,text width=13mm] (t3) at (179,0) {\textbf{T3 sondear}\\conteos a\\2 turnos};
  % carril de callejón
  \node[tool,c3,text width=21mm] (t4) at (55,-15) {\textbf{T4 explicar}\\\textbf{callejón}\\causa + 12 candidatas};
  \node[agent,c3a,text width=21mm] (a3) at (91,-15) {\textbf{A3 Diagnosticador}\\elige a qué turno\\saltar};
  \node[tool,c3,text width=17mm] (t5) at (122,-15) {\textbf{T5 lecciones}\\memoria verbal};
  \node[tool,cp,text width=15mm] (t7) at (1,-14) {\textbf{T7 completar}\\relajado (D2)};
  \node[tool,cp,text width=15mm] (t6) at (1,-23.9) {\textbf{T6 verificador}\\HC1--HC6};
  % leyenda
  \node[tool,minimum height=4.5mm,minimum width=6mm] (lt) at (150,-17) {T};
  \node[anchor=west] at (154.5,-17) {\textbf{Herramienta} (código)};
  \node[agent,minimum height=5mm,minimum width=6mm,rounded corners=2mm] (la) at (150,-23.5) {A};
  \node[anchor=west] at (154.5,-23.5) {\textbf{Agente} (LLM)};
  % carril normal: aristas (ciclo 1, azul)
  \begin{scope}[cuno]
  \draw[arr] (in) -- (orq);
  \draw[arr] (orq) -- node[lbl,above=0.4pt]{turno $k$} (t1);
  \draw[arr] (t1) -- node[lbl,above=0.4pt]{$\geq$2 opciones:\\menú barajado} (a1);
  \draw[arr] (a1.north east) -- node[lbl,above left=-0.5pt]{propuesta} (t2.west);
  \draw[line width=0.5pt] (t2.north) -- (120,19);
  \draw[arr] (120,19) -| (19,0 |- orq.north);
  \node[lbl] at (84,19) {sin alertas, o \textbf{acepta}: fijar la opción y pasar a $k{+}1$};
  \node[cyc,fill=white,inner sep=0.5pt] at (36,19) {CICLO 1: turno a turno};
  \draw[arr] (t1.north) |- (27,8.2) -- (27,0 |- orq.north);
  \node[lbl] at (41,8.2) {1 opción: fijar sin LLM};
  \end{scope}
  % revisión: aristas (ciclo 2, naranjo)
  \begin{scope}[cdos]
  \draw[arr] (t2.east) -- node[lbl,above right=-0.5pt]{alertas +\\números} (a2.north west);
  \draw[arr] (a2.west) -- node[lbl,below=0.8pt]{\textbf{rechaza} + motivo} node[cyc,above=0.8pt]{CICLO 2 ($\leq$2/turno)} (a1.east);
  \draw[arr] ([yshift=1.3mm]a2.east) -- node[lbl,above=0.4pt]{sondear ($\leq$1)} ([yshift=1.3mm]t3.west);
  \draw[arr] ([yshift=-1.3mm]t3.west) -- node[lbl,below=0.4pt]{conteos} ([yshift=-1.3mm]a2.east);
  \draw[line width=0.5pt] (a2.north) |- (120,19);
  \draw[line width=0.5pt] (a1.north) -- (91,19);
  \node[lbl,anchor=west] at (91.6,12.6) {A1 mantiene\\o 3.ª propuesta};
  \end{scope}
  % carril de callejón: aristas (ciclo 3, rojo)
  \begin{scope}[ctres]
  \draw[arr] (t1) -- node[lbl,right=0.5pt]{0 opciones:\\callejón} (t4);
  \draw[arr] (t4) -- node[lbl,above=0.4pt]{explicación} (a3);
  \draw[arr] (a3) -- node[lbl,above=0.4pt]{lección} (t5);
  \draw[arr] (t5.north west) -- node[lbl,left=0.3pt,pos=0.5]{$\leq$3\\lecciones} (a1.south east);
  \draw[arr] (t5.north east) -- node[lbl,right=0.3pt,pos=0.5]{$\leq$3\\lecciones} (a2.south west);
  \draw[arr] (a3.south) |- (29,-23.5) -- (29,0 |- orq.south);
  \node[lbl] at (60,-23.5) {\textbf{salta} al turno $j$ y prohíbe ahí a un médico};
  \node[cyc,anchor=west] at (93,-23.5) {CICLO 3: saltos ($\leq$30)};
  \end{scope}
  % parada (verde)
  \begin{scope}[cpar]
  \draw[arr] (orq.south west) -- node[lbl,left=0.2pt,pos=0.4]{presupuesto\\agotado} (t7.north east);
  \draw[arr] (t7) -- node[lbl,right=0.3pt]{no cert.} (t6);
  \draw[arr] (17,0 |- orq.south) |- (t6.east);
  \node[lbl,anchor=west] at (17.6,-17.5) {\textbf{PARADA}\\$k{=}21$};
  \end{scope}
\end{tikzpicture}
\end{center}
\vspace{-15pt}

\raggedcolumns
\begin{multicols}{2}

\section{1. Continuidad con D1 y D2, y desviaciones declaradas}
\textbf{Qué se mantiene de D1 y D2:}
\begin{itemize}[label=$\blacktriangleright$]
    \item la tarea: horario de 10 médicos en 21 turnos (7 días $\times$ 3), salida JSON;
    \item el juez: \texttt{verifier.py}; correcto $=$ no viola ninguna regla dura (HC1--HC6);
    \item el modelo: \textbf{Llama-3.2-3B-Instruct} (4-bit, greedy).
    \item las entradas: 30 principales y 30 de estrés (de D2), más 30 difíciles.
\end{itemize}
\textbf{La falla que se ataca:} la elección voraz de D1 (F1): gastar temprano las horas y el descanso de los médicos deja el fin de semana imposible. En D2 siguió \emph{dentro} de la elección: una opción legal podía condenar la semana 2--12 turnos antes del callejón, y retroceder de a un turno no llega a esa \emph{decisión culpable}. D3 vuelve directo a ella.\par
\textbf{Cómo se construye sobre D2:} D3 reutiliza el código de D2 (\texttt{stepwise\_solver.py}, sin cambios) para filtrar las opciones legales y para completar el horario cuando se acaba el presupuesto. Prueba: si se le quita lo nuevo y elige al azar como D2, entrega los mismos 90 horarios que D2 (\texttt{test\_continuity.py}).\par
\textbf{Desviaciones respecto de D2:}
\begin{itemize}[label=$\blacktriangleright$]
    \item \textbf{Conjunto difícil:} D2 ya resolvía 30/30 del principal: se agregaron 30 instancias con capacidad justa, elegidas por su estructura, no por resultados.
    \item \textbf{Respuestas de 1 token:} el modelo responde siempre un número, el de mayor probabilidad; equivale a la decodificación restringida de D2.
    \item \textbf{Presupuesto:} 30 saltos, como los 30 retrocesos de D2; como un salto deshace hasta 12 decisiones, también se compara con D2 al azar con 300.
    \item \textbf{Ajuste de prompts:} con 20 instancias de calibración aparte, que no están entre las 90 evaluadas.
    \item \textbf{Correcciones:} antes de la corrida final, una revisión del código arregló errores del sistema y de la medición.
\end{itemize}

\section{2. Las herramientas T1--T7 (\texttt{src/tools.py})}
Código sin LLM: calculan números sobre el horario parcial para los agentes.\par
{\tabfont
\begin{tabularx}{\linewidth}{@{}L{19mm}Y@{}}
\toprule
\textbf{T1} opciones viables & Opciones para $k$ que cumplen HC1--HC6 y el chequeo hacia adelante (lógica de D2), sin las prohibidas por saltos. \\
\addlinespace[1pt]
\textbf{T2} perfilar & Alertas vs.\ la mejor alternativa: R1 más especialistas sobre la cuota, R2 deja menos de la mitad de opciones, R3 más turnos sin margen. \\
\addlinespace[1pt]
\textbf{T3} sondear & Continuaciones viables a 2 turnos (propuesta y alternativa). \\
\addlinespace[1pt]
\textbf{T4} explicar callejón & Requisito que ya no se cubre, médicos que faltan y por qué, y las últimas 12 decisiones como candidatas a rehacer. \\
\addlinespace[1pt]
\textbf{T5} lecciones & Memoria verbal: dónde fue el callejón, a qué turno se volvió y a quién se prohibió. A1 y A2 leen las 3 últimas. \\
\addlinespace[1pt]
\textbf{T6} verificador & \texttt{verifier.py}, el juez de D1 y D2: violaciones de HC1--HC6 del horario final (correcto $=$ 0). \\
\addlinespace[1pt]
\textbf{T7} completar relajado & Sin presupuesto, completa relajando horas, descanso y luego bloqueos; elige A1 (lógica de D2). \\
\bottomrule
\end{tabularx}\par}
{\notefont \textbf{HC1} descanso $\geq$16\,h, \textbf{HC2} 1 turno/día, \textbf{HC3} dotación, \textbf{HC4} cuota de especialidad, \textbf{HC5} horas, \textbf{HC6} bloqueos.\par}

\section{3. Los tres agentes (prompts versionados en \texttt{prompts/})}
Mismo modelo, prompts distintos; cada uno responde un solo número.\par
{\tabfont
\begin{tabularx}{\linewidth}{@{}L{17.5mm}YL{16.5mm}Y@{}}
\toprule
\textbf{Agente} & \textbf{Recibe} & \textbf{Responde} & \textbf{Luego} \\
\midrule
\textbf{A1}\newline Asignador & estado, menú barajado (T1), lecciones (T5) y motivo del rechazo & n.º de opción;\newline \texttt{0} $=$ mantener & T2 la perfila; sin alertas, se fija \\
\addlinespace[1.5pt]
\textbf{A2}\newline Crítico & propuesta frente a la mejor alternativa, con números de T2 y T3 & \texttt{1} acepta\newline \texttt{2} rechaza\newline \texttt{3} sondea & acepta: se fija; rechaza: vuelve a A1 (máx.\ 2) \\
\addlinespace[1.5pt]
\textbf{A3}\newline Diagnosticador & de T4: causa, médicos en conflicto y últimas 12 decisiones & n.º de la decisión a rehacer ($j$) & deshace $j..k$, prohíbe un médico en $j$, lección a T5 (máx.\ 30) \\
\bottomrule
\end{tabularx}\par}
\textbf{Herramientas:} A2 puede llamar a T3 (sondear); A1 y A3 no llaman herramientas: reciben lo que el coordinador obtiene de T1, T4 y T5.\par
\textbf{Por qué se divide así:} A1 mira a lo ancho; A2 revisa a fondo una propuesta (a un 3B le cuesta menos verificar que generar); A3 mira hacia atrás.

\section{4. Cómo leer el grafo de agentes}
Rectángulos: herramientas (código); cajas redondeadas: agentes (LLM). Cada color es un ciclo. Todo lo dirige el \textbf{coordinador de turnos} (gris), código que guarda el horario a medio armar y el turno $k$, y decide a quién le toca. Se llena de a un turno, del lunes AM ($k{=}1$) al domingo PM ($k{=}21$):
\begin{itemize}[label=$\blacktriangleright$]
    \item[\sw{cuno}] \textcolor{cuno}{\textbf{Ciclo 1, turno a turno:}} T1 da las opciones viables de $k$. Con 1 se fija sin LLM; con $\geq$2, A1 elige y T2 la perfila. Sin alertas, se fija y se pasa a $k{+}1$.
    \item[\sw{cdos}] \textcolor{cdos}{\textbf{Ciclo 2, revisión:}} con alertas decide A2 (puede sondear con T3). Acepta: se fija. Rechaza: vuelve a A1 con el motivo (máx.\ 2; la 3.ª se fija).
    \item[\sw{ctres}] \textcolor{ctres}{\textbf{Ciclo 3, callejón:}} con 0 opciones, T4 lo explica y A3 elige a qué turno $j$ volver: se deshace $j..k$, se prohíbe un médico en $j$ y T5 anota la lección.
    \item[\sw{cpar}] \textcolor{cpar}{\textbf{Parada:}} 21 turnos fijados $=$ certificado. Con los 30 saltos gastados, T7 completa relajado (no certificado). T6 juzga ambos.
\end{itemize}

\section{5. Técnicas de prompting}
\textbf{5a. Qué técnicas usa cada agente} (categorías de \emph{The Prompt Report})\par
{\tabfont
\begin{tabularx}{\linewidth}{@{}L{25mm}cccY@{}}
\toprule
\textbf{Técnica} (\emph{categoría}) & \textbf{A1} & \textbf{A2} & \textbf{A3} & \textbf{Falla que ataca} \\
\midrule
Role prompting\newline\emph{zero-shot} & $\checkmark$ & $\checkmark$ & $\checkmark$ & F3 de D1: en salidas largas el modelo inventaba IDs de médicos y rompía el JSON. \\
\addlinespace[1pt]
Salida 1 token\newline\emph{answer engineering} & $\checkmark$ & $\checkmark$ & $\checkmark$ & F3: responde solo un número, así que no puede inventar IDs ni romper el JSON. \\
\addlinespace[1pt]
Few-shot\newline\emph{few-shot (ICL)} & $\checkmark$ & $\checkmark$ & $\checkmark$ & Elegir mal la opción o el salto; los ejemplos muestran casos ya resueltos. \\
\addlinespace[1pt]
CRITIC\newline\emph{agentes: tool use} & & $\checkmark$ & & El chequeo de D2 deja pasar opciones legales que ya condenan la semana. \\
\addlinespace[1pt]
Self-Refine\newline\emph{self-criticism} & $\checkmark$ & & & La misma: A1 rehace su propuesta con el motivo del rechazo de A2. \\
\addlinespace[1pt]
ReAct\newline\emph{agentes: observación} & & & $\checkmark$ & F1: la decisión culpable está 2--12 turnos atrás y retroceder 1, como D2, no llega. \\
\addlinespace[1pt]
Reflexion\newline\emph{agentes: observación} & $\checkmark$ & $\checkmark$ & $\checkmark$ & Repetir, después de un salto, el mismo gasto que causó el callejón. \\
\addlinespace[1pt]
Ensamble\newline\emph{ensembling} & $\checkmark$ & & & Sesgo de posición: D2 elegía los primeros candidatos el 33\,\% de las veces (azar: 17\,\%). \\
\bottomrule
\end{tabularx}\par}
\textbf{5b. Qué pasa al quitar cada técnica} (mismas 90 instancias)\par
{\fontsize{6.3}{7.2}\selectfont\setlength{\tabcolsep}{2.2pt}
\begin{tabularx}{\linewidth}{@{}L{17mm}>{\centering\arraybackslash}p{12.8mm}>{\centering\arraybackslash}p{10.5mm}>{\centering\arraybackslash}p{12.5mm}Y@{}}
\toprule
\textbf{Sistema} & \textbf{Correctos} & \textbf{Sin\newline callejón} & \textbf{Llamadas al LLM} & \textbf{Qué aporta y por qué} \\
\midrule
Línea base & 0 & -- & 1 & \\
Solución D2 & 80 & 42 & 26 & \\
\textbf{D3 completo} & \textbf{90} & \textbf{78} & \textbf{77} & \\
\midrule
$-$ ReAct (sin A3) & 82 ($-$8) & 78 (=) & 66 & \textbf{Clave:} sin A3 se retrocede 1 turno, como D2: vuelve F1. \\
\addlinespace[1pt]
$-$ Few-shot & 87 ($-$3) & 67 ($-$11) & 112 & \textbf{Ayuda:} sin ejemplos, +11 callejones y +45\,\% llamadas. \\
\addlinespace[1pt]
$-$ CRITIC+ Self-Refine & 89 ($-$1) & 57 ($-$21) & 63 & \textbf{Ayuda:} A2 frena malas elecciones; sin él, +21 callejones. \\
\addlinespace[1pt]
$-$ Reflexion & 90 (=) & 81 (+3) & 86 & \textbf{Mejora saltos:} sin lecciones, sirve el 70\,\% de los saltos (no el 100\,\%). \\
\addlinespace[1pt]
$-$ Ensamble & 90 (=) & 76 ($-$2) & 48 & \textbf{No aporta:} casi igual, pero cuesta 60\,\% más llamadas. \\
\bottomrule
\end{tabularx}\par}
{\notefont Rol y salida de 1 token son la base: sin ellas, es la línea base (0/90).\par}

\section{6. Resultados (mismas 90 instancias, mismo juez)}
Horarios correctos según \texttt{verifier.py}. Abajo, \emph{controles}: el mismo sistema con el LLM reemplazado por azar (promedio de 5 corridas) o por reglas fijas.\par
{\tabfont
\begin{tabular}{@{}lcccc@{}}
\toprule
\textbf{Sistema} & \textbf{Princ.} & \textbf{Estrés} & \textbf{Difícil} & \textbf{Total} \\
\midrule
Línea base: prompt directo (D1) & 0 & 0 & 0 & 0/90 \\
Solución D2 & 30 & 28 & 22 & 80/90 \\
\textbf{Sistema de agentes (D3)} & \textbf{30} & \textbf{30} & \textbf{30} & \textbf{90/90} \\
\midrule
\multicolumn{5}{@{}l}{\emph{Controles: el LLM se reemplaza por azar o por reglas}} \\
C1. D2 eligiendo al azar & 29.8 & 27.4 & 24.6 & 81.8 \\
C2. D2 al azar con 300 retrocesos & 29.8 & 28.2 & 27.0 & 85.0 \\
C3. D3 con agentes al azar & 30 & 30 & 30 & 90.0 \\
C4. D3 con reglas fijas & 29 & 25 & 20 & 74 \\
\bottomrule
\end{tabular}\par}
\begin{itemize}[label=$\blacktriangleright$]
    \item D3 resuelve los 90 horarios; D2, 80; la línea base, 0. D3 gana en estrés y difícil (30 y 30, contra 28 y 22).
    \item D2 sin LLM da casi lo mismo (C1: 81.8), y ni con 10 veces más retrocesos alcanza a D3 (C2: 85).
    \item D3 con agentes al azar también llega a 90 (C3): el total lo da la estructura de D3; el LLM aporta menos callejones (5b).
    \item Con reglas fijas (C4) baja a 74: saltan a la decisión equivocada.
\end{itemize}
\section{7. Límites y reproducción}
\textbf{Límites} (lo que este trabajo no demuestra): el 90/90 no se debe al LLM: con agentes al azar (C3) también se llega. El LLM aporta menos callejones y mejores saltos (A3 arregló los 27 callejones; al azar, 69\,\%; \texttt{decision\_eval.py}). Con 90 instancias, diferencias de 1 a 3 horarios pueden ser azar.\par
\textbf{Reproducción:} los pasos del video están en \texttt{video/}; las corridas, trazas y tablas, en \texttt{results/}; y el paso a paso, en el README.

\end{multicols}
\end{document}
